\documentclass{article}
%\usepackage[utf8]{inputenc} %no lo requiero, se está compilando con lualatex
\usepackage[spanish]{babel}
\usepackage{listings}
\usepackage{graphicx}
\graphicspath{ {images/} }
\usepackage{cite}
\usepackage{float}
\usepackage{enumerate}

\begin{document}

%***************************PORTADA***********************************

\begin{titlepage}
    \begin{center}
        \vspace*{1cm}
            
        \Huge
        \textbf{Nociones de la memoria del computador
}
            
        \vspace{0.5cm}
        \LARGE
        %Subtítulo
            
        \vspace{1.5cm}
            
        \textbf{Diego Alejandro Londoño Jiménez\\ \emph{diegoa.londono@udea.edu.co} }
            
        \vfill
            
        \vspace{0.8cm}
            
        \Large
        Despartamento de Ingeniería Electrónica y Telecomunicaciones\\
        Universidad de Antioquia\\
        Medellín\\
        Septiembre de 2020\\
        
            
    \end{center}
\end{titlepage}


%***************************FIN PORTADA***********************************




\tableofcontents

\newpage

\section{Preguntas del taller}\label{preguntas}
    \begin{enumerate}[a)]
    \item Defina que es la memoria del computador.
    \item Mencione los tipos de memoria que conoce y haga una pequeña descripción de cada tipo.
    \item Describa la manera como se gestiona la memoria en un computador
    \item ¿Qué hace que una memoria sea más rápida que otra? ¿Por qué esto es importante?
    \end{enumerate}

\section{Respuestas del taller} \label{respuestas}
    Esta sección es para ver qué pasa con los comandos que definen texto.

\subsection{Citación}
    Vamos a citar por ejemplo un artículo de \textbf{Albert Einstein} \cite{einstein}.
    También es posible citar libros \cite{dirac} o documentos en línea \cite{knuthwebsite}.\\\\
    Revisar en la última sección el formato de las referencias en IEEE.

\subsection{EJEMPLO}


En la sección \ref{imagenes}, se presentará como añadir ilustraciones al texto.

\section{Inclusión de imágenes} \label{imagenes}

En la Figura (\ref{fig:cpplogo}), se presenta el logo de C++ contenido en la carpeta images.

\begin{figure}[H]
\includegraphics[width=4cm]{cpplogo.png}
\centering
\caption{Logo de C++}
\label{fig:cpplogo}
\end{figure}

Las secciones (\ref{intro}), (\ref{contenido}) y (\ref{imagenes}) dependen del estilo del documento.

\bibliographystyle{IEEEtran}
\bibliography{references}

\end{document}
